\section{Optical distortion and summation at one common physical roof}
\label{sec:pt-roof}
The synchronized momentum sum in \eqref{eq:rc-synchronized} is not
itself a sum at one physical roof: a fixed roof generally selects a
different momentum radius on every critical branch.  We now supply an
explicit comparison mechanism on the uniform physical tubes just
constructed.  It uses logarithmic optical distortion before summation,
not an exchange of the roof supremum and the sum.  The only global
root input is the retained physical kernel Theorem~\ref{thm:mc-kernel},
with its native growth dependencies left unchanged.

Set $\sigma=47/59$, $r=q_0^{1/4}$ and $\beta=\sigma/r<1$.
The constants $J_*,C_c,C_F$ retain the exact meanings given in
Section~\ref{sec:el-localization}.  Define
\begin{align}
 g&=R_-/2=9/40, &U&=320/9,\notag\\
 D_*&=\frac{8C_c}{1-\beta}
       +\frac{4aJ_*}{1-\sigma}+\frac{4C_F}{g},
 &C_{\rm dist}&=\exp(D_*\delta_*/2).\label{eq:pt-dist-constants}
\end{align}
No constant in this display depends on collision count, labels or
an unproved lower bound for an interior incidence.

\begin{lemma}[Relative distortion of the actual inverse momentum weight]
\label{lem:pt-distortion}
On any convex physical endpoint domain with
$c_i\ge(\varepsilon/2)r^{d_i}$ and $c_0,c_m\ge1/2$, put
$G=D_q^2F$, $e=G_{0m}$ and
\[
 u(q)=\frac{R|e(q)|}{\det G(q)}
     =|\partial_{\alpha_0}p_m|^{-1}.
\]
For $0<\varepsilon\le1/4$,
\begin{equation}\label{eq:pt-log-distortion}
 |D_q\log u[\eta]|\le D_*\varepsilon^{-1}\|\eta\|_\infty.
\end{equation}
Consequently $u(q)\le e^{D_*\|q-q_x\|_\infty/\varepsilon}u(q_x)$
whenever the connecting segment is in that domain.
\end{lemma}
\begin{proof}
The true optical formula is
$|e|=R^2c_0c_m/\mathfrak b_m$.  Positivity and separate
multiaffinity of $\mathfrak b_m$, proved in
Lemma~\ref{lem:el-logden}, give
\[
 |D\log|e||\le\sum_{i=0}^m\frac{|Dc_i|}{c_i}
                  +\sum_{i=0}^{m-1}\frac{|D\tau_i|}{\tau_i}.
\]
Each incidence occurs once in this bound, including the two endpoints;
there is no new product of inverse incidences.  With
$s_i=\sigma^i+\sigma^{m-i}\le2\sigma^{d_i}$ and at most two
indices at every depth,
\[
 \sum_{i=0}^m s_i r^{-d_i}\le\frac4{1-\beta},\qquad
 \sum_{i=0}^{m-1}(s_i+s_{i+1})\le\frac4{1-\sigma}.
\]
Proposition~\ref{prop:el-observables} therefore bounds the incidence
sum by $8C_c/[\varepsilon(1-\beta)]$ and the flight sum by
$4aJ_* /(1-\sigma)$, times $\|\eta\|_\infty$.
The physical Schur complement has $G\succeq gI$.  By
Theorem~\ref{thm:el-third}, $\|DG[\eta]\|_{\rm op}\le
2C_F\|\eta\|_\infty$, so
\[
 |D\log\det G[\eta]|
 =|\operatorname{tr}(G^{-1}DG[\eta])|
 \le4C_Fg^{-1}\|\eta\|_\infty.
\]
Subtract these logarithmic derivatives in $\log u$ and use
$\varepsilon\le1$ to obtain \eqref{eq:pt-log-distortion}.
Integration on the actual physical segment proves the last assertion.
The identity for $u$ is the unchanged generating-map Jacobian of
Theorem~\ref{thm:rc-coarea}; no radius or section factor is absorbed in it.
\end{proof}

For $0<\varepsilon\le1/4$ let $\mathscr Z_{\varepsilon,m,R}$ be
all distinct regular physical critical roots $p_0=p_m=0$ satisfying
\eqref{eq:pt-guarded-center}.  Lattice translation copies and angular
seams are identified, so each root is a unique physical initial point.
These are actual reversing critical orbits by the retained theorem,
not critical points of an unrelated surrogate action.  Their roots are
isolated on the initial momentum fibre since
$\partial_{\alpha_0}p_m\ne0$; hence the family is countable.
Put
\begin{equation}\label{eq:pt-momentum-radius}
 s_\varepsilon=\tfrac12\delta_*\varepsilon,
 \qquad
 \rho_\varepsilon=\min\{1/4,\,g s_\varepsilon/(2R_+)\},
 \qquad K_{1/2}=\frac{100}{47}\sqrt{1-1/4}.
\end{equation}
Both positive radii are fixed before taking any collision limsup.

\begin{lemma}[A common physical momentum disk at every guarded critical root]
\label{lem:pt-momentum-disk}
For every $x\in\mathscr Z_{\varepsilon,m,R}$ there is an actual
regular inverse $q_x(\zeta)$ for $|\zeta|<\rho_\varepsilon$,
where $\zeta=(p_0,p_m)$ and $q_x(0)$ is the original endpoint pair.
It is contained in the physical endpoint ball
$|q-q_x(0)|_2<s_\varepsilon$ and obeys
\[
 |q_x(\zeta)-q_x(0)|_2\le R_+|\zeta|/g,
 \qquad u_x(\zeta)\le C_{\rm dist}u_x(0).
\]
This is a common domain for all collision words in the stated class,
not a separate radius chosen after $m$ is fixed.
\end{lemma}
\begin{proof}
The Euclidean ball of radius $s_\varepsilon$ lies strictly inside
the $\ell^1$ endpoint tube of Theorem~\ref{thm:pt-physical}.
At a critical center both endpoint tangential components vanish.
The signed tangential estimate in that proof and
$C_{\rm inc}\delta_*\le1/2$ imply
$|p_0|,|p_m|\le\varepsilon/2$ on the entire ball; in particular
$c_0,c_m\ge1/2$ and $G\succeq gI$ there.  The domain and every
segment used next are genuine physical domains, established before
using the Hessian.

Given $|\zeta|<\rho_\varepsilon$, minimize
$F(q)-R(D\zeta)\cdot(q-q_x(0))$ on the closed endpoint ball,
where $D=\operatorname{diag}(-1,1)$.  On its boundary the radial
first variation is at least
$g s_\varepsilon^2-R_+|\zeta|s_\varepsilon>0$.
A boundary minimizer is impossible, because moving radially inward
would strictly decrease the function.  Strong convexity gives one
interior minimizer, satisfying $\nabla F=RD\zeta$, and hence the
correct endpoint momenta.  The ordinary inverse theorem gives its
smooth dependence, while strong convexity gives
$g|q-q_x(0)|_2\le R_+|\zeta|$.  The optical mixed entry is
nonzero, so this is also a regular local parametrization of the
original collision source.  Lemma~\ref{lem:pt-distortion}, integrated
on the physical radial segment in endpoint space, gives the stated
relative bound.  In it
$\|q-q_x(0)\|_\infty\le s_\varepsilon$ and
$s_\varepsilon/\varepsilon=\delta_*/2$.
\end{proof}

\begin{theorem}[One-copy critical-tube sum at a common physical roof]
\label{thm:pt-common-roof}
For fixed $R,m,\varepsilon$ take the union of the actual momentum
charts in Lemma~\ref{lem:pt-momentum-disk}.  Give this union its
original collision measure, counting each physical initial point once.
Retain any Borel selector $0\le w\le1$, in particular the full
$\Pi_{n,k,m,R}$, both section tests, the original guard and any
original physical-roof interval.  Let $Q^w_{\varepsilon,m,R}(t)$
be the density of its pushforward by the original total flight time.
Then
\begin{equation}\label{eq:pt-common-roof}
 \|Q^w_{\varepsilon,m,R}\|_{L^\infty(dt)}
       \le \frac{U C_{\rm dist}C_{\rm mom}}{2c_*}.
\end{equation}
This bound is uniform in $\varepsilon\in(0,1/4]$, all $m\ge1$,
all original radii and exact labels.  It makes the complete positive
chart sum before taking the physical-roof essential supremum.
It is unscaled: it does not assert an $m^{-2}$ gain.
\end{theorem}
\begin{proof}
Enumerate the distinct critical roots.  Each chart is a Borel image
of its open momentum disk; the endpoint-to-source map is locally
invertible since $e\ne0$.  It is injective on the convex endpoint
ball: at fixed $\alpha_0$, $\partial_{\alpha_m}p_0=-e/R$ has a
fixed nonzero sign along the connected vertical slice.  Thus a
physical source point is represented at most once in a given chart.
If charts overlap, assign the physical point to the first chart in
the enumeration.  These Borel assignments partition the unchanged
union; they remove only coordinate multiplicity, never an actual
collision or a point of the union.  Multiply $w$ by the assignment
indicator in each chart, denoting the result by $w_x$.

Let $a_x=L(x)$ and write $\zeta=r e_\theta$.  The exact radial
coarea in Theorem~\ref{thm:rc-coarea} gives, for $r>0$,
\[
 \partial_r L_x(re_\theta)=r h_x(re_\theta),\qquad
 U^{-1}\le h_x\le K_{1/2}^{-1},\qquad
 b_x(re_\theta)=\frac{u_x(re_\theta)}{4\pi c_*h_x(re_\theta)}.
\]
These formulas hold on the full common disk by the preceding lemma.
For fixed $t,\theta$, strict radial monotonicity gives at most one
$r_x(t,\theta)\in(0,\rho_\varepsilon)$ per center.  These radii
are generally different for different centers; we do not identify them.
The contribution is zero when no such radius exists.  Otherwise
\[
 w_x b_x(r_x(t,\theta)e_\theta)
       \le \frac{U C_{\rm dist}}{4\pi c_*}\,u_x(0).
\]
The retained complete regular-root theorem at the single common
momentum pair $(0,0)$ gives
$\sum_{x\in\mathscr Z_{\varepsilon,m,R}}u_x(0)\le C_{\rm mom}$.
Here the transfer from different momentum points to $(0,0)$ is paid
by the preceding distortion estimate, not inferred from the
synchronized kernel alone.

Apply the coarea identity first to bounded nonnegative roof tests
and finitely many charts, retaining the Borel selectors in the
integrands.  The upper bound just obtained is summable independently
of $t$ and $\theta$.  Monotone convergence gives the full nonnegative
source sum.  Integrating in $\theta\in[0,2\pi)$ gives precisely
\eqref{eq:pt-common-roof}; the polar origin has zero collision area
in each regular chart and creates no atom.  This area statement does
not set a separate one-sided pairing trace to zero.  At no step is a
sum of individual roof suprema bounded by a synchronized value without
a comparison factor.
\end{proof}

\begin{corollary}[Uniform band coverage and a localized critical-roof envelope]
\label{cor:pt-envelope}
Every chart above reaches every physical roof
\[
 0<t-a_x<\rho_\varepsilon^2/(2U)
\]
on every momentum ray.  Moreover the density in the theorem satisfies
for almost every $t$ the sharper bound
\begin{equation}\label{eq:pt-envelope}
 Q^w_{\varepsilon,m,R}(t)
 \le\frac{U C_{\rm dist}}{2c_*}
 \sum_{\substack{x\in\mathscr Z_{\varepsilon,m,R}\\
       t-\rho_\varepsilon^2/(2K_{1/2})\le a_x\le t}}u_x(0).
\end{equation}
The endpoints $a_x$ are the actual physical critical roofs, not clock
corrections.  No section visit count is transferred from a tube point
to its center unless section stability is proved separately.
\end{corollary}
\begin{proof}
Integrating $r/U\le\partial_rL_x\le r/K_{1/2}$ from zero gives
$r^2/(2U)\le L_x(re_\theta)-a_x\le r^2/(2K_{1/2})$.
Continuity and strict monotonicity prove coverage; the upper inequality
restricts the contributing centers at any fixed physical roof.
Insert that restriction into the nonnegative sum in the theorem.
The original Borel exact-record selector is still evaluated on each
actual orbit; the envelope simply dominates it by one.
\end{proof}

\subsection{The remaining normalized source is not eliminated}
The new analytic domain, physical guard tube and common-roof summation
remove two specific gaps of the earlier critical argument: no
collision-dependent continuation radius is chosen, and different
branch momentum radii are now compared by a proved multiplicative
optical bound.  They neither assert that all critical centers satisfy
a common guard margin nor estimate every original singular, noncritical
or paired-boundary source.  The latter measures need not be restrictions
of the bulk collision measure used here.

Even on the covered union, \eqref{eq:pt-common-roof} is an $O(1)$
roof-height bound, not an $O(m^{-2})$ bound.  It cannot be multiplied
by $m^2$ and declared negligible.  The envelope
\eqref{eq:pt-envelope} isolates genuine physical roof locations but
provides no unproved critical-roof nonconcentration rate or transfer
of return labels across a section cut.  The missing original limit
remains exactly \eqref{eq:v112-remaining}, not a renamed substitute.
To close it requires a normalized complete physical-source estimate or
sufficient original nongrazing capacity, including the uncovered
incidence-only low-load source, unvisited and visited boundary terms,
and the retained spectral, arithmetic and conditioned-denominator
interfaces.  Its original widths stay fixed before $m\to\infty$
and $B\to\infty$ remains outermost.  No limit has been interchanged
in the new finite-width, all-collision theorem.
